\documentclass[10pt,a4paper]{amsart}
\usepackage{amsmath,amssymb,amsthm,graphicx}
% For commutative diagrams we use Xy-pic
\usepackage[all]{xy}

% We use 2cell for 2-commutative diagrams.
\xyoption{2cell}
% Theorem environments.
%
\theoremstyle{plain}
\newtheorem{theorem}[subsection]{Theorem}
\newtheorem{proposition}[subsection]{Proposition}
\newtheorem{lemma}[subsection]{Lemma}

\theoremstyle{definition}
\newtheorem{definition}[subsection]{Definition}
\newtheorem{example}[subsection]{Example}
\newtheorem{exercise}[subsection]{Exercise}
\newtheorem{situation}[subsection]{Situation}

\theoremstyle{remark}
\newtheorem{remark}[subsection]{Remark}
\newtheorem{remarks}[subsection]{Remarks}

\numberwithin{equation}{subsection}

% Macros
%
\def\lim{\mathop{\mathrm{lim}}\nolimits}
\def\colim{\mathop{\mathrm{colim}}\nolimits}
\def\Spec{\mathop{\mathrm{Spec}}}
\def\Hom{\mathop{\mathrm{Hom}}\nolimits}
\def\Ext{\mathop{\mathrm{Ext}}\nolimits}
\def\SheafHom{\mathop{\mathcal{H}\!\mathit{om}}\nolimits}
\def\SheafExt{\mathop{\mathcal{E}\!\mathit{xt}}\nolimits}
\def\Sch{\mathit{Sch}}
\def\Mor{\mathop{\mathrm{Mor}}\nolimits}
\def\Ob{\mathop{\mathrm{Ob}}\nolimits}
\def\Sh{\mathop{\mathit{Sh}}\nolimits}
\def\NL{\mathop{N\!L}\nolimits}
\def\CH{\mathop{\mathrm{CH}}\nolimits}
\def\proetale{{pro\text{-}\acute{e}tale}}
\def\etale{{\acute{e}tale}}
\def\QCoh{\mathit{QCoh}}
\def\Ker{\mathop{\mathrm{Ker}}}
\def\Im{\mathop{\mathrm{Im}}}
\def\Coker{\mathop{\mathrm{Coker}}}
\def\Coim{\mathop{\mathrm{Coim}}}

% Boxtimes
%
\DeclareMathSymbol{\boxtimes}{\mathbin}{AMSa}{"02}

%
% Macros for moduli stacks/spaces
%
\def\QCohstack{\mathcal{QC}\!\mathit{oh}}
\def\Cohstack{\mathcal{C}\!\mathit{oh}}
\def\Spacesstack{\mathcal{S}\!\mathit{paces}}
\def\Quotfunctor{\mathrm{Quot}}
\def\Hilbfunctor{\mathrm{Hilb}}
\def\Curvesstack{\mathcal{C}\!\mathit{urves}}
\def\Polarizedstack{\mathcal{P}\!\mathit{olarized}}
\def\Complexesstack{\mathcal{C}\!\mathit{omplexes}}
% \Pic is the operator that assigns to X its picard group, usage \Pic(X)
% \Picardstack_{X/B} denotes the Picard stack of X over B
% \Picardfunctor_{X/B} denotes the Picard functor of X over B
\def\Pic{\mathop{\mathrm{Pic}}\nolimits}
\def\Picardstack{\mathcal{P}\!\mathit{ic}}
\def\Picardfunctor{\mathrm{Pic}}
\def\Deformationcategory{\mathcal{D}\!\mathit{ef}}
\usepackage{cleveref}
\crefname{originalreference}{original reference}{original references}
\Crefname{originalreference}{Original reference}{Original references}
\makeatletter
\newlabel{lemma-cohomological-dimension-proper}{{095U}{source}{}{sourcepointer-1}{}}
\newlabel{lemma-cohomological-dimension-proper@cref}{{[originalreference][0][]095U}{[1][0][]source}{}{}{}}
\makeatother
\crefname{originalreference}{original reference}{original references}
\Crefname{originalreference}{Original reference}{Original references}
\makeatletter
\newlabel{lemma-finite-cd-mod-n-direct-sums}{{0F11}{source}{}{sourcepointer-2}{}}
\newlabel{lemma-finite-cd-mod-n-direct-sums@cref}{{[originalreference][0][]0F11}{[1][0][]source}{}{}{}}
\makeatother
\crefname{originallemma}{original lemma}{original lemmas}
\Crefname{originallemma}{Original lemma}{Original lemmas}
\makeatletter
\newlabel{morphisms-lemma-morphism-finite-type-bounded-dimension}{{0A3V}{source}{}{sourcepointer-3}{}}
\newlabel{morphisms-lemma-morphism-finite-type-bounded-dimension@cref}{{[originallemma][0][]0A3V}{[1][0][]source}{}{}{}}
\makeatother
\title{The pushforward along a proper morphism commutes with direct sums for coefficients in Z/nZ}
\author{The Stacks Project Authors}
\date{stacks-project:etale-cohomology.tex:local-snapshot-20261011; GNU FDL 1.2 or later}
\begin{document}
\maketitle
\noindent\textit{Editorial note.} All delivered mathematical text is the author's own, copied from the pinned Stacks Project chapter file etale-cohomology.tex; no statement or proof step was rewritten, reordered or removed.\par\medskip
\noindent\textit{Editorial note.} The counted proof is the author's complete argument: the proof is written out in full: the reduction to a quasi-compact base, the fibre dimension bound, the finiteness of the cohomological dimension, and the application of the commutation lemma.\par\medskip
\noindent\textit{Editorial note.} The label lemma-proper-mod-n-direct-sums names the counted statement, and the proof environment immediately following it in the source is delivered verbatim as its proof.\par\medskip
\noindent\textit{Editorial note.} The counted statement is carried by the author's own lemma environment in the chapter source.\par\medskip
\noindent\textit{Editorial note.} The statement and proof are copied from the pinned chapter file; the three cited inputs are kept as external pointers with their canonical Stacks tags.\par\medskip
\noindent\textit{Allowed standard prerequisites (editorial).}\begin{itemize}
\item Boundedness of the fibre dimension of a morphism of finite type. Morphisms, the lemma bounding the dimension of the fibres of a morphism of schemes. Supplies the boundedness of the fibre dimension used to get finite cohomological dimension.
\item Finiteness of the cohomological dimension of a proper morphism. Same chapter, the lemma bounding the cohomological dimension of a proper morphism in terms of the fibre dimension. Supplies cd(f) < infinity.
\item Commutation of the pushforward with direct sums under finite cohomological dimension. Same chapter, the lemma stating that the pushforward with Z/nZ coefficients commutes with direct sums when the cohomological dimension is finite. Supplies the conclusion.
\end{itemize}\medskip
\ifdefined\hypertarget\hypertarget{sourcepointer-1}{}\fi
\noindent\textit{Source reference.} Original reference 095U: sources/stacks-etale-cohomology/pinned-original-source.tex, source lines 18786--18786. Prerequisite pointer invoked at proof line(s) 20158; it names a standard fact of the same chapter that the counted argument uses, and every step of the counted proof built on it is retained verbatim in the delivered body.\par\medskip
\ifdefined\hypertarget\hypertarget{sourcepointer-2}{}\fi
\noindent\textit{Source reference.} Original reference 0F11: sources/stacks-etale-cohomology/pinned-original-source.tex, source lines 20074--20074. Prerequisite pointer invoked at proof line(s) 20160; it names a standard fact of the same chapter that the counted argument uses, and every step of the counted proof built on it is retained verbatim in the delivered body.\par\medskip
\ifdefined\hypertarget\hypertarget{sourcepointer-3}{}\fi
\noindent\textit{Source reference.} Original lemma 0A3V: sources/stacks-chow/tags.txt, source lines 12753--12753. Prerequisite pointer invoked at proof line(s) 20155; it names a standard fact of a sibling chapter that the counted argument uses, and every step of the counted proof built on it is retained verbatim in the delivered body.\par\medskip
\begin{lemma}
\label{lemma-proper-mod-n-direct-sums}
Let $f : X \to Y$ be a proper morphism of schemes. Let $n \geq 1$
be an integer. Then the functor
$$
Rf_* :
D(X_\etale, \mathbf{Z}/n\mathbf{Z})
\longrightarrow
D(Y_\etale, \mathbf{Z}/n\mathbf{Z})
$$
commutes with direct sums.
\end{lemma}

\begin{proof}
It is enough to prove this when $Y$ is quasi-compact. By
Morphisms, Lemma \ref{morphisms-lemma-morphism-finite-type-bounded-dimension}
we see that the dimension of the fibres of
$f : X \to Y$ is bounded.
Thus Lemma \ref{lemma-cohomological-dimension-proper}
implies that $\text{cd}(f) < \infty$. Hence the
result by Lemma \ref{lemma-finite-cd-mod-n-direct-sums}.
\end{proof}

\end{document}
